% A.42: registered after independent full root and DE audits, both PASS.
% This adds the coefficient-prime-equals-residue-characteristic case to A.41.
\subsection{The residue-characteristic coefficient prime}

\begin{statement}{A.42}{Proved}{The compact-type rational-path criterion holds at every prime}
Under the geometric hypotheses of A.41, its fixed-subgroup
conclusion also holds for \(\ell=p\). Consequently for
every prime \(\ell\), every cusp packet at \(P\), and
every indicated open derived subgroup, the full fixed group
of the actual geometric pro-\(\ell\) fundamental group
is exactly \(I_{P,\ell}\).

The crystalline argument uses the connected global algebraic
monodromy group only after a change of fibre. It does not
regard individual global Galois matrices as morphisms of
local crystalline representations, and it does not apply
Fontaine's monodromy to an individual etale vector.
\end{statement}

\begin{proof}
The assertion for \(\ell\ne p\) is A.41. Put
\(\ell=p\), and make the finite global field restriction
used there, so the marked model is split and semistable,
its graph has the stated rational \(P\)--\(Q\) path,
the proper homology image has connected Zariski closure,
and \(\Gamma=\overline{[G_L,G_L]}\) lies in the
specified open derived subgroup. Write \(F=L_v\) for
the local field and \(F_0\) for its maximal unramified
subfield. Put
\[
 H=\bigl(\pi_1(X_{\bar L})^{(p)}\bigr)^{s(\Gamma)}.
\]
Let \(\mathfrak h\) be the Lie algebra of its actual
unipotent closure, formed before taking finite weight
quotients, exactly as in A.38.

\paragraph{Crystalline leading subobjects killed by local monodromy.}
The general comparison proved in the first two paragraphs
of A.31 applies to this marked semistable curve and every
packet at its selected cusp. It makes every finite rational
unipotent Lie quotient semistable over \(F\), and identifies
its \(D_{\rm st}\) realization with the log-crystalline
one, whose monodromy is the actual edge-twist residue.
Those arguments were stated for an arbitrary residue prime
and marked semistable curve; their comparison and packet
proofs do not use the later special \(S_3\) models.
In particular arbitrary rationalized cusp Kummer twists
remain semistable and contribute a cusp-inner residue.

Use the weight filtration of A.41: proper homology has
weight one and cusp generators have weight two. Its
intrinsic second step is the kernel of the map to proper
abelian homology; subsequent steps are generated by the
bracket and by the weight-two cusp images. The comparison
is compatible with this filtration by its functoriality
for proper compactification, brackets, and the cusp maps.
Thus no crystalline splitting of ordinary open homology
is required.

Every finite quotient of \(\mathfrak h\) is a globally
stable subrepresentation of a finite ambient Lie quotient.
Its \(G_L\)-action factors through \(G_L^{\rm ab}\),
because \(\Gamma\) fixes the actual subgroup pointwise.
It is semistable locally by closure of semistability under
subquotients, and crystalline by the lemma in A.31 that a semistable
representation with abelian image is crystalline. Hence
Fontaine's \(N\) vanishes on its whole \(D_{\rm st}\)
realization.

The proper Jacobian has good reduction, since the special
curve is of compact type. Thus
\(V=H_1(\overline X_{\bar L},\mathbf Q_p)\) is
crystalline at \(v\). The independent boundary space
\(T\) is a sum of \(\mathbf Q_p(1)\)'s and is also
crystalline. Every finite weighted graded piece of the
ambient free Lie algebra, and every subobject
\[
                          W_d=\operatorname{gr}_d\mathfrak h,
\]
is consequently crystalline. Here and below the free Lie
algebra uses \(V\) in weight one and \(T\) in weight two.

In the log-crystalline realization write
\[
 V_{\rm c}=D_{\rm cris}(V)=\bigoplus_i V_{i,\rm c},
 \quad T_{\rm c}=D_{\rm cris}(T),\qquad
 \Omega_{\rm c}=\sum_i\Omega_{i,\rm c}.              \tag{A.42.1}
\]
The blocks come from the positive-genus components of the
special compact-type curve and are orthogonal symplectic.
Choose a basis of the one-dimensional crystalline Tate
fibre to view their polarization classes as bivectors.
The edge formula of A.31, with its node-thickness
coefficients, gives exactly (A.41.4)--(A.41.6) over
\(F_0\), including the same positive distances \(r_i\).
The actual tangent has \(N(t_P)=0\). A possible inner
weight-one discrepancy centralizing \(t_P\) is zero,
and a weight-two discrepancy belongs to
\(F_0t_P\). Thus the first possible weighted residue
is the genuine derivation
\[
                  N_2=D_{\mathrm{graph}}+
                                      \lambda\operatorname{ad}(t_P),
                         \qquad\lambda\in F_0,       \tag{A.42.2}
\]
up to a common nonzero normalization. There is no
weight-zero residue because the weighted generators
are crystalline, and no weight-one outer residue for
the separating graph. The displayed inner comparison
also excludes a weight-one based residue.

Exactness of \(D_{\rm st}\), applied to a sufficiently
large finite quotient of \(\mathfrak h\) and its
induced weight filtration, now gives
\[
                         D_{\rm cris}(W_d)\subset\ker N_2.
                                                               \tag{A.42.3}
\]
More explicitly, lift a vector of the left side to the
weight-\(d\) part of that finite \(D_{\rm st}\) quotient.
Its monodromy is zero; its component in weight \(d+2\)
is \(N_2\) of the leading vector because all smaller
weight shifts vanish. Surjectivity of the finite
unipotent-closure transitions, as proved in A.38, ensures
that the full \(W_d\) is represented in this procedure.

\paragraph{Transport the global group through a genuine fibre functor.}
Let \(\mathcal C\) be the full tensor subcategory of
global \(\mathbf Q_p\)-representations generated by
\(V\) and \(\mathbf Q_p(1)\), closed under subquotients
and duals. Its usual fibre functor \(\omega\) has group
\(G\), the connected global algebraic monodromy group.
The proper polarization makes \(G\) a subgroup of
\(\operatorname{GSp}(V)\), with multiplier the Tate
character. Every object of \(\mathcal C\) is locally
crystalline. Therefore
\[
                    \eta=D_{\rm cris}|_{\mathcal C}
                               :\mathcal C\longrightarrow\operatorname{Vec}_{F_0}
                                                               \tag{A.42.4}
\]
is a faithful exact tensor fibre functor, by
\cite[Theorem 5.2.1 and Proposition 9.1.6]{fghBrinonConrad}.
Its tensor automorphism group \(G_{\rm c}\) is an
inner form of \(G_{F_0}\): the fibre-comparison torsor
\(\operatorname{Isom}^{\otimes}(\omega_{F_0},\eta)\)
is faithfully flat by
\cite[Theorem 3.2, p.31]{fghDeligneMilneTC}.
In particular \(G_{\rm c}\) is connected. It acts on
\(V_{\rm c}\) by symplectic similitudes and on
\(T_{\rm c}\) by its Tate character \(\chi_{\rm c}\).
Each \(W_d\) is an object of \(\mathcal C\), so its
crystalline fibre in (A.42.3) is \(G_{\rm c}\)-stable.

This statement can be checked directly with the period
comparison, avoiding any ambiguity about individual
matrices. Let \(E=\operatorname{Frac}(B_{\rm cris})\).
The functorial crystalline comparison gives tensor
isomorphisms
\[
 c:E\otimes_{F_0}D_{\rm cris}(V)
                       \xrightarrow{\sim}E\otimes_{\mathbf Q_p}V,
                                                               \tag{A.42.5}
\]
and likewise for \(T\), all their tensor/Lie
constructions, and every subobject \(W_d\). Thus
\(c\) identifies the crystalline leading subspace
with \(E\otimes W_d\). Conjugating \(G_E\) by
\(c\) gives the connected comparison group acting
on that same crystalline ambient space. This is not
an assertion that a global matrix is defined on the
\(F_0\)-fibre. The map \(N_2\) in (A.42.3) is
extended \(E\)-linearly for this linear-algebra
calculation; one does not differentiate period
coefficients or conjugate Fontaine's differential
operator onto etale vectors.

\paragraph{The diagonal argument and descent of scalar characters.}
Apply the all-weight calculation of A.41 over \(F_0\)
to (A.42.2), with
\(U_{\rm c}=\langle\Omega_{i,\rm c}\rangle\).
It gives
\[
 D_{\rm cris}(W_d)\subset
                  \mathbb L(T_{\rm c}\oplus U_{\rm c})_d.
                                                               \tag{A.42.6}
\]
Because its left side is \(G_{\rm c}\)-stable, the
exact tensor-intersection argument puts it in
\[
 \mathbb L(T_{\rm c}\oplus U_{\rm c,*})_d,
           \qquad U_{\rm c,*}=\bigcap_{g\in G_{\rm c}}gU_{\rm c}.
                                                               \tag{A.42.7}
\]
Compute these intersections and connected orbits after extending
\(F_0\) to an algebraic closure; all ensuing scalar identities
descend faithfully. The connected diagonal-projector lemma of A.41 applies
to \(G_{\rm c}\) and the local orthogonal component
blocks. It proves that \(G_{\rm c}\) acts on
\(U_{\rm c,*}\) by \(\chi_{\rm c}\). All generators
on the right of (A.42.7) therefore have weight two and
that character. It follows that odd \(W_d\)'s vanish
and \(G_{\rm c}\) acts on \(D_{\rm cris}(W_{2n})\)
by the scalar \(\chi_{\rm c}^n\).

A scalar-character identity is preserved by tensor
change of fibre and faithfully flat descent. Concretely,
trivialize the comparison torsor, or use (A.42.5) over
\(E\). Conjugation of a scalar matrix leaves it
unchanged. The preceding identity therefore becomes
that the global \(G\)-representation on
\(W_{2n}\) is \(\chi_p^n\) after a faithful scalar
extension, and hence already over \(\mathbf Q_p\).
This is precisely the Tate-grading hypothesis of A.38
for the closure of the actual subgroup \(H\). That
criterion proves \(H=I_{P,p}\). Enlarging back to the
specified open derived subgroup, which fixes the same
inertia, proves the assertion.

The algebraic kernel and intersection certificates of
\texttt{code/a\_compact\_type\_fixed.py} apply unchanged:
all their identities are over characteristic-zero fields.
The additional step here is the functorial fibre change
and its scalar-character descent, proved above.
\end{proof}
